\section{Introduction}

A product of complex manifolds has a tangent bundle split into the
pullbacks of the tangent bundles of its factors. A converse must keep
track of more than the abstract bundles. For a specified decomposition
\[
T_X=E_1\oplus E_2,
\]
the embeddings of \(E_1\) and \(E_2\) in \(T_X\) are part of the data,
and a compatible product must identify these subbundles with its factor
directions. A necessary local condition is \emph{integrability}:
the holomorphic local sections of each \(E_i\) must be closed under
the Lie bracket.

We prove that rational connectedness forces this local condition.
Here and below, ``two general points'' means all pairs in some nonempty
Zariski open subset of \(X\times X\). A smooth connected projective complex
manifold is rationally connected if two general points lie on the image
of a morphism \(\PP^1\to X\).

\begin{theorem}[Automatic integrability]\label[theorem]{thm:main}
Let \(X\) be a smooth connected projective complex manifold of dimension
at least two. Suppose that two general points of \(X\) lie on the image
of a morphism \(\PP^1\to X\). For every specified holomorphic decomposition
\[
T_X=E_1\oplus E_2
\]
into subbundles of positive rank, both \(E_1\) and \(E_2\) are integrable.
\end{theorem}

The theorem has no further rank or positivity hypothesis. It proves the
smooth rationally connected assertion formulated in
\cite[Conjecture~1.5]{Hor26}. H\"oring's earlier product theorem then
answers the global question for the same specified splitting.

\Needspace{7\baselineskip}
\begin{corollary}[Compatible product]\label[corollary]{cor:product}
Under the hypotheses of \cref{thm:main}, there are smooth projective
manifolds \(X_1,X_2\) and an isomorphism \(X\simeq X_1\times X_2\)
under which
\[
E_i=\pr_i^*T_{X_i}\qquad (i=1,2).
\]
\end{corollary}

\begin{proof}
H\"oring's \cite[Theorem~1.4]{Hor07} applies to a smooth connected
rationally connected projective complex manifold with a specified
splitting of its tangent bundle and at least one integrable summand.
By \cref{thm:main}, these hypotheses hold. His theorem gives an
isomorphism with a product that identifies the specified summands with
the factor tangent bundles. Each factor is smooth. It is projective
because its realization as a slice of the product is a closed complex
submanifold of \(X\), and a projective embedding of \(X\) restricts
to an embedding of that slice.
\end{proof}

\subsection*{Context and scope}

Beauville studied how a tangent-bundle splitting on a compact K\"ahler
manifold can lead to a product of its universal cover \cite{Bea00}.
An abstract product need not make a given summand integrable:
H\"oring's \cite[Example~2.14]{Hor07} gives a nonintegrable summand on
a product of an abelian surface with \(\PP^1\). Rational connectedness
rules out that example and supplies the abundance of rational curves
used in our proof.

Several earlier results establish compatible products under positivity
or rank assumptions. Campana and Peternell treated Fano manifolds with
a two-summand splitting when one summand has rank one or two, and hence
every two-summand splitting on a Fano manifold of dimension at most five
\cite[Theorems~3.5--3.6 and Corollary~3.7]{CP02}. H\"oring's
\cite[Theorem~1.4]{Hor07} showed that, for a rationally connected
projective manifold, integrability of one summand already gives a
compatible product. He then isolated the remaining question by
conjecturing that at least one summand is integrable
\cite[Conjecture~1.2]{Hor08}. His later low-rank result applies to a
uniruled projective manifold whose tangent bundle splits into summands
all of rank at most two, and proves that at least one of them is
integrable \cite[Lemma~4.21]{Hor08}.

The recent Fano-type theorem \cite[Theorem~1.1]{Hor26} establishes
integrability with algebraic leaves for split tangent sheaves on normal
\(\mathbb Q\)-factorial projective klt varieties of Fano type. In the
smooth Fano case it yields an actual compatible product
\cite[Corollary~1.3]{Hor26}. The present theorem removes this positivity
condition for smooth rationally connected manifolds. Smoothness matters:
\cite[Example~4.3]{Hor26} gives a singular rationally connected example
with a nonintegrable summand. Section~\ref{sec:consequences} applies the
smooth theorem to general fibers of the rationally connected quotient,
supplying the integrability premise in H\"oring's structure theorems
for uniruled compact K\"ahler manifolds.

\subsection*{Proof strategy}

The construction that drives the proof does not require rational
connectedness. Given any holomorphic map \(g:\PP^1\to X\) into a smooth
projective manifold with a split tangent bundle, we construct
\[
F:\PP^1\times\PP^1\longrightarrow X,\qquad F(s,s)=g(s),
\]
with the first factor direction tangent to \(E_1\) and the second
tangent to \(E_2\). We call \(F\) a \emph{tangent box}. The name records
these tangencies; the map may have critical points and need not give
local coordinates.

Write \(P_i:T_X\to E_i\subset T_X\) for the holomorphic projectors
defined by the splitting. The diagonal first carries a formal solution. In coordinates
\(u=s\), \(v=t-s\), the equation
\(\partial_v\widehat F=P_2(\widehat F)\partial_u\widehat F\)
recursively separates the derivative of \(g\) into its two summands.
The positive normal bundle of the diagonal then bounds the number of
global formal sections of a twisted line bundle. Comparing that bound
with pairs of polynomial sections forces the formal projective
coordinate ratios to be rational. A local cancellation makes these
ratios holomorphic near a diagonal patch. Finally, on a resolution of
the rational map, the two tangency identities force every exceptional
curve to be contracted. The resulting map is holomorphic on the entire
product. The section-count argument is given directly for holomorphic
formal series; the algebraic formal-function literature provides
broader context \cite{HM68,Bad96}.

The splitting also provides the projected-bracket partial connection
\(\nabla^i_\xi e=P_i[\xi,e]\) on \(E_i\) in directions of the
complementary summand. On a tangent box it gives transport along the
fibers of \(\pr_i\), which vary the complementary factor, and hence
\[
F^*E_i\simeq\pr_i^*(g^*E_i),
\]
normalized to be the identity on the diagonal. Thus, on a factor curve,
the corresponding summand retains the bundle \(g^*E_i\), while the
complementary summand is trivial.

Rational connectedness enters only in the last step. A two-point family
of rational curves supplies sections of \(g^*E_i\) that vanish at one
marked point and span the summand at another. Applying the bracket tensor
\(\bigwedge^2E_i\to E_j\) to two such sections on a factor curve gives
a section of a trivial bundle on \(\PP^1\). It is constant and vanishes
at the first mark, so it vanishes at the second. The spanning property
then kills the bracket tensor on an open set of \(X\), and hence
everywhere. The use of rational curves to annihilate this tensor has a
direct antecedent in \cite[Section~2.B and Lemma~3.18]{Hor07}; the tangent
box supplies the trivial complementary target in arbitrary rank.

\subsection*{Organization}

\Cref{sec:boxes} constructs tangent boxes through the formal Cauchy
problem, rational reconstruction, and removal of indeterminacy.
\Cref{sec:transport} establishes the bundle identifications, and
\cref{sec:vanishing} proves \cref{thm:main}. The rational-quotient
consequence follows in \cref{sec:consequences}. Appendix~\ref{app:surface}
proves the required point-blowup resolution for a surface, and
Appendix~\ref{app:variation} derives the two-point family directly from
rational connectedness.
